% Options for packages loaded elsewhere
\PassOptionsToPackage{unicode}{hyperref}
\PassOptionsToPackage{hyphens}{url}
%
\documentclass[
]{article}
\usepackage{amsmath,amssymb}
\usepackage{iftex}
\ifPDFTeX
  \usepackage[T1]{fontenc}
  \usepackage[utf8]{inputenc}
  \usepackage{textcomp} % provide euro and other symbols
\else % if luatex or xetex
  \usepackage{unicode-math} % this also loads fontspec
  \defaultfontfeatures{Scale=MatchLowercase}
  \defaultfontfeatures[\rmfamily]{Ligatures=TeX,Scale=1}
\fi
\usepackage{lmodern}
\ifPDFTeX\else
  % xetex/luatex font selection
\fi
% Use upquote if available, for straight quotes in verbatim environments
\IfFileExists{upquote.sty}{\usepackage{upquote}}{}
\IfFileExists{microtype.sty}{% use microtype if available
  \usepackage[]{microtype}
  \UseMicrotypeSet[protrusion]{basicmath} % disable protrusion for tt fonts
}{}
\makeatletter
\@ifundefined{KOMAClassName}{% if non-KOMA class
  \IfFileExists{parskip.sty}{%
    \usepackage{parskip}
  }{% else
    \setlength{\parindent}{0pt}
    \setlength{\parskip}{6pt plus 2pt minus 1pt}}
}{% if KOMA class
  \KOMAoptions{parskip=half}}
\makeatother
\usepackage{xcolor}
\usepackage{longtable,booktabs,array}
\usepackage{calc} % for calculating minipage widths
% Correct order of tables after \paragraph or \subparagraph
\usepackage{etoolbox}
\makeatletter
\patchcmd\longtable{\par}{\if@noskipsec\mbox{}\fi\par}{}{}
\makeatother
% Allow footnotes in longtable head/foot
\IfFileExists{footnotehyper.sty}{\usepackage{footnotehyper}}{\usepackage{footnote}}
\makesavenoteenv{longtable}
\usepackage{graphicx}
\makeatletter
\def\maxwidth{\ifdim\Gin@nat@width>\linewidth\linewidth\else\Gin@nat@width\fi}
\def\maxheight{\ifdim\Gin@nat@height>\textheight\textheight\else\Gin@nat@height\fi}
\makeatother
% Scale images if necessary, so that they will not overflow the page
% margins by default, and it is still possible to overwrite the defaults
% using explicit options in \includegraphics[width, height, ...]{}
\setkeys{Gin}{width=\maxwidth,height=\maxheight,keepaspectratio}
\newcommand{\researchfigure}[1]{%
  \IfFileExists{#1}{\includegraphics{#1}}{%
    \fbox{\parbox{0.92\linewidth}{\centering Figure file not present in this checkout:\\
      \texttt{\detokenize{#1}}}}}}
% Set default figure placement to htbp
\makeatletter
\def\fps@figure{htbp}
\makeatother
\setlength{\emergencystretch}{3em} % prevent overfull lines
\providecommand{\tightlist}{%
  \setlength{\itemsep}{0pt}\setlength{\parskip}{0pt}}
\setcounter{secnumdepth}{-\maxdimen} % remove section numbering
\ifLuaTeX
  \usepackage{selnolig}  % disable illegal ligatures
\fi
\usepackage[]{natbib}
\bibliographystyle{plainnat}
\IfFileExists{bookmark.sty}{\usepackage{bookmark}}{\usepackage{hyperref}}
\IfFileExists{xurl.sty}{\usepackage{xurl}}{} % add URL line breaks if available
\urlstyle{same}
\hypersetup{
  pdftitle={Privacy--Utility Auditing of DP-SGD for Machine-Learning-Based Network Intrusion Detection},
  pdfauthor={Author names and affiliations to be confirmed},
  hidelinks,
  pdfcreator={LaTeX via pandoc}}

\title{Privacy--Utility Auditing of DP-SGD for Machine-Learning-Based
Network Intrusion Detection}
\author{Author names and affiliations to be confirmed}
\date{}

\begin{document}
\maketitle

\begin{abstract}

Machine-learning-based network intrusion detection systems can expose
information about their training records, but privacy protection is
useful only if detection performance remains operationally meaningful.
We evaluate this tension for a binary multilayer-perceptron intrusion
detector trained with and without formally accounted differentially
private stochastic gradient descent (DP-SGD). The primary study uses a
locked NSL-KDD partition, five target-training seeds, condition-matched
shadow models, score-only and label-aware black-box membership-inference
attacks (MIAs), validation-only F2 threshold selection, and an untouched
KDDTest+ utility test. A supplementary single-seed experiment repeats
the main conditions on the official UNSW-NB15 training/testing
partitions. On NSL-KDD, the non-private, epsilon-about-4, and
epsilon-about-2 conditions achieve mean F1 scores of 0.8097, 0.8017, and
0.7946. Their false-positive rates are 0.0561, 0.0876, and 0.0838, and
their average precisions are 0.9375, 0.8907, and 0.8906. The small
epsilon-about-4 Recall increase is uncertain across seeds; paired
analyses instead show a higher false-positive burden and lower average
precision for both private conditions. Overall MIA AUCs remain close to
0.5 for every condition, and all paired DP-minus-non-private AUC
intervals cross zero. On UNSW-NB15, private training similarly preserves
selected-threshold F1 and near-perfect Recall, but the operating points
have approximately 42\% false-positive rates and reduced average
precision. These results do not show a DP-induced reduction in measured
membership leakage, because the evaluated non-private model is already
at an empirical attack floor. The formal DP guarantee therefore remains
the principal privacy evidence, conditional on fixed non-private
preprocessing, while the MIA results delimit what the tested attacks
could measure. The study contributes a reproducible evaluation design
and a bounded negative result: under this protocol, DP-SGD changes the
operational utility profile without yielding a detectable improvement in
population-level membership inference.
\end{abstract}

\textbf{Keywords:} differential privacy; DP-SGD; intrusion detection;
membership inference; NSL-KDD; UNSW-NB15; privacy--utility tradeoff

\hypertarget{introduction}{%
\subsection{1. Introduction}\label{introduction}}

Machine-learning-based intrusion detection systems (IDSs) learn
statistical patterns from network traffic and can detect behavior not
covered by a fixed signature database. Their training data may, however,
encode sensitive facts about hosts, users, communications, and recorded
attacks. A model can leak information about those records even when the
original dataset is not released. Membership inference formalizes one
such risk: given a candidate record and access to a trained model, an
adversary tries to decide whether the record was part of the training
set \citep{shokri2017membership}.

Differential privacy (DP) offers a distributional guarantee that limits
the influence of one record on a randomized algorithm's output
\citep{dwork2006calibrating}. For neural networks, DP-SGD applies
per-example gradient clipping, Gaussian noise, subsampling, and
privacy-loss composition during optimization \citep{abadi2016deep}. The
guarantee is attractive because it does not depend on defeating one
chosen empirical attack. Yet DP-SGD can degrade predictive utility, and
an IDS has asymmetric operational costs: a false negative misses
malicious activity, while a false positive consumes analyst attention
and may make a detector unusable. Accuracy alone cannot represent this
tradeoff.

Empirical MIAs remain useful as audits of observable leakage, but
interpreting them requires care. Shadow-model attacks can learn
behavioral differences between members and non-members
\citep{shokri2017membership}, metric-based attacks can outperform poorly
configured learned attacks \citep{song2021systematic}, and average-case
metrics can hide leakage concentrated at low false-positive rates
\citep{carlini2022membership}. Recent work further shows that
population-average or defense-unaware evaluations can materially
underestimate the most vulnerable records \citep{aerni2024evaluations}.
Consequently, a near-chance MIA result is evidence only about the stated
attacker and evaluation population; it is not proof that a model is
private.

Existing privacy-preserving IDS studies establish that DP can be
incorporated into network-traffic classification, but they address
different mechanisms and questions. Markovic et al.~combine federated
random forests with a tree-specific exponential mechanism and report
accuracy/F1 on four IDS datasets, without a membership audit
\citep{markovic2024random}. Liu et al.~introduce a feature-level
Laplace-noise framework evaluated against membership and adversarial
attacks on USTC-TFC2016 and CICIDS2017 \citep{liu2025hierarchical}.
Neither study is a centralized, multi-seed DP-SGD evaluation with paired
MIA uncertainty and IDS-specific operating metrics.

This paper asks:

\begin{itemize}
\tightlist
\item
  \textbf{RQ1:} How does formally accounted DP-SGD affect Recall,
  false-negative rate (FNR), F1, false-positive rate (FPR), and average
  precision for a fixed binary IDS pipeline?
\item
  \textbf{RQ2:} Under predeclared score-only and label-aware black-box
  MIAs, is measured membership leakage lower for the private models than
  for a matched non-private model?
\item
  \textbf{RQ3:} Does a constrained, single-seed UNSW-NB15 replication
  show the same broad privacy--utility pattern as the five-seed NSL-KDD
  study?
\end{itemize}

The contribution is not a new privacy mechanism and not a claim that DP
is new to IDS. Instead, we contribute: (1) a reproducible protocol
combining formal accounting, IDS operating-point selection,
condition-matched shadow attacks, and paired uncertainty; (2) a
five-seed NSL-KDD analysis that reports false-alarm burden and ranking
quality alongside Recall/FNR; (3) a bounded single-seed external check
on UNSW-NB15; and (4) an explicit negative finding that near-chance
baseline MIA leaves insufficient empirical headroom to demonstrate
leakage reduction under the tested attacks.

\hypertarget{background-and-related-work}{%
\subsection{2. Background and related
work}\label{background-and-related-work}}

\hypertarget{differential-privacy-and-dp-sgd}{%
\subsubsection{2.1 Differential privacy and
DP-SGD}\label{differential-privacy-and-dp-sgd}}

A randomized mechanism $M$ is $(\varepsilon,\delta)$-differentially
private if, for all adjacent datasets $D,D'$ differing in one record and
all measurable output events $S$,

\[
\Pr[M(D)\in S] \leq e^{\varepsilon}\Pr[M(D')\in S]+\delta.
\]

Smaller $\varepsilon$ gives a tighter worst-case bound for fixed
$\delta$, but epsilon values are meaningful only with the adjacency
relation, training procedure, sampling assumption, and accountant
stated. DP-SGD clips each per-example gradient to a maximum norm and
adds calibrated Gaussian noise before the optimizer update
\citep{abadi2016deep}. We use Opacus \citep{yousefpour2021opacus} and
its privacy-random-variable (PRV) accountant, which is based on
numerical composition of privacy-loss random variables
\citep{gopi2021numerical}.

Our formal statement covers the stochastic optimization stage.
Categorical vocabularies, scaling parameters, the model architecture,
and hyperparameters are fixed outside the private optimizer. We
therefore do not claim end-to-end privacy for raw preprocessing. We also
record \texttt{secure\_mode=False}; the runs support research
evaluation, not a production-strength randomness claim.

\hypertarget{membership-inference-as-attack-relative-evidence}{%
\subsubsection{2.2 Membership inference as attack-relative
evidence}\label{membership-inference-as-attack-relative-evidence}}

The original black-box shadow-model formulation trains auxiliary models
to generate member and non-member examples for an attack classifier
\citep{shokri2017membership}. Later studies show that loss, confidence,
correctness, entropy, and per-class calibration can be competitive or
stronger signals \citep{song2021systematic}. LiRA demonstrates that
likelihood-ratio modeling and low-FPR evaluation can uncover behavior
hidden by aggregate accuracy or AUC \citep{carlini2022membership}.
Canary-based auditing can also construct empirical lower bounds on a
claimed DP guarantee, but requires a different protocol from ordinary
population MIAs \citep{steinke2023privacy}.

Our study evaluates ordinary shadow-calibrated, population-level MIAs.
We report ROC-AUC, advantage, balanced accuracy, and TPR at 1\% and 5\%
FPR. We do not run LiRA, RMIA, canary-based worst-case auditing,
label-only perturbation attacks, or a white-box adaptive adversary. The
result therefore answers RQ2 only under the stated attacks; it does not
empirically certify the claimed epsilon or rule out leakage from
individual vulnerable samples \citep{aerni2024evaluations}.

\hypertarget{privacy-preserving-and-adversarially-robust-ids}{%
\subsubsection{2.3 Privacy-preserving and adversarially robust
IDS}\label{privacy-preserving-and-adversarially-robust-ids}}

Markovic et al.~evaluate differentially private random forests in a
horizontal federated-learning framework on KDD, NSL-KDD, UNSW-NB15, and
CIC-IDS-2017. Their exponential-mechanism trees use epsilon values 0.1,
0.5, 1, and 5, and the evaluation reports accuracy and F1 for attack
detection and classification \citep{markovic2024random}. The work
demonstrates a DP--IDS utility tradeoff with public code, but does not
evaluate membership inference and is not directly comparable to
centralized DP-SGD.

Liu et al.~explicitly combine membership-inference and
adversarial-attack defense. Their HierarchicalDP framework ranks
structured input features and applies different Laplace noise to
privacy- or category-sensitive features. It is evaluated on USTC-TFC2016
and CICIDS2017 with MPE and EMI membership attacks and NIFGSM/APGD
evasion attacks \citep{liu2025hierarchical}. This is the closest
integrated privacy/security neighbor, but it is a feature-perturbation
method, not private gradient training, and it reports
membership-inference success rates rather than our multi-seed paired AUC
analysis.

Adversarial evasion is related but distinct from membership privacy.
Tafreshian and Zhang use a genetic algorithm intended to respect feature
mutability, interdependence, and protocol constraints, then combine
adversarial training, balancing, feature engineering, ensemble learning,
and tuning on NSL-KDD and UNSW-NB15 \citep{tafreshian2025defensive}.
Sharma and Chen systematically apply white-box and black-box attacks to
nine NSL-KDD classifiers, but note that practical feature masking
remains future work \citep{sharma2024systematic}. These studies motivate
precise threat models; they do not provide evidence that our DP-SGD
models are robust to evasion. Adversarial robustness is outside the
present experimental scope.

\hypertarget{method}{%
\subsection{3. Method}\label{method}}

\hypertarget{datasets-and-locked-partitions}{%
\subsubsection{3.1 Datasets and locked
partitions}\label{datasets-and-locked-partitions}}

NSL-KDD was proposed to reduce redundancy-related problems in KDD'99 and
provides manageable training and test files
\citep{tavallaee2009detailed}. We use all 125,973 rows of
\texttt{KDDTrain+} as the development source. With seed 42, it is
stratified by binary label and coarse attack family into 88,181
target-training records (70\%), 12,597 target-validation records (10\%),
and 25,195 shadow-pool records (20\%). \texttt{KDDTest+} contains 22,544
records and is used only for final IDS utility. It is never used to
train or calibrate an attacker or select an IDS threshold.

UNSW-NB15 contains modern synthetic/real traffic generated in the UNSW
Cyber Range and is distributed with official training and testing
partitions \citep{moustafa2015unsw}. For the supplementary study, the
175,341-row training partition is divided into 122,738 target-training,
17,534 target-validation, and 35,069 shadow-pool records. The complete
82,332-row official testing partition is used only for utility.
Experiment 09 uses one target-training seed and is descriptive external
evidence, not an equally powered replication.

Both datasets are reduced to binary Normal-versus-Attack classification.
Coarse NSL-KDD families are retained for descriptive subgroup analysis,
but subgroup findings are not primary because of small rare-family
samples and multiple comparisons.

\hypertarget{preprocessing-and-model}{%
\subsubsection{3.2 Preprocessing and
model}\label{preprocessing-and-model}}

For NSL-KDD, the 41 traffic features are used and the supplied
difficulty field is excluded. \texttt{protocol\_type}, \texttt{service},
and \texttt{flag} are one-hot encoded with unknown-category handling;
numeric features are min--max scaled. The preprocessor is fitted only on
target-train and then applied to validation and test records, producing
122 inputs. The same structural procedure is fitted within each shadow
split. UNSW-NB15 uses its frozen Experiment 09 schema; identifiers and
labels are not features, categorical variables are one-hot encoded,
numeric variables are scaled, and the fitted preprocessor is learned
from the corresponding training data only.

The target classifier is a feed-forward MLP: input--64--32--1, ReLU
activations, and one output logit. There is no batch normalization.
Non-private and private conditions use Adam for 30 epochs with batch
size 256, learning rate 0.001, weight decay 0.0001, and binary
cross-entropy with logits. Architecture, optimizer, epochs, split, and
threshold policy are held fixed across conditions.

\hypertarget{private-training-and-accounting}{%
\subsubsection{3.3 Private training and
accounting}\label{private-training-and-accounting}}

The primary analysis compares non-private training with target privacy
budgets 4 and 2. Private runs use maximum gradient norm 1.0, Poisson
sampling, and PRV accounting. For NSL-KDD,
$\delta=1/88181=1.1340311\times10^{-5}$. The resulting epsilon
values are 3.998267 and 1.999038, with noise multipliers 0.6817627 and
0.8813477. Each condition is trained with target seeds 42, 52, 62, 72,
and 82. Epsilon about 8 was evaluated only in the earlier single-seed
sweep and is retained as context, not included in the five-seed primary
comparison.

The UNSW-NB15 private targets reach epsilon 3.995481 and 1.995670 at
$\delta=8.147436\times10^{-6}$. Its target and shadow models follow
the accepted Experiment 09 configuration. All privacy-accounting inputs,
warnings, target/shadow configurations, and completion gates are stored
in the repository evidence.

\hypertarget{ids-operating-point-evaluation}{%
\subsubsection{3.4 IDS operating-point
evaluation}\label{ids-operating-point-evaluation}}

For each trained target, candidate thresholds are evaluated only on
target-validation. The chosen threshold maximizes F2, weighting Recall
more heavily than precision. The selected threshold is then applied once
to the external test partition. We report Recall, FNR, F1, FPR,
precision, and average precision (PR-AUC). Average precision is
threshold-independent and helps distinguish a change in ranking quality
from a change caused by the chosen operating threshold.
Default-threshold results are retained as a diagnostic, particularly for
UNSW-NB15.

\hypertarget{membership-inference-protocol}{%
\subsubsection{3.5 Membership-inference
protocol}\label{membership-inference-protocol}}

Five condition-matched shadow models are trained for each target
condition. Shadows 101, 202, 303, and 404 create the attacker-training
data; shadow 505 is reserved for attack calibration. No target output is
used to select an attacker, fit its parameters, or calibrate its
operating threshold. Private shadows target the same epsilon and delta
as the corresponding target while deriving a shadow-specific noise
multiplier.

The target evaluation population is identical across conditions: 12,597
sampled members from target-train and all 12,597 target-validation
non-members. The score-only threat model observes the predicted attack
probability; its frozen primary attacker is logistic regression. The
label-aware audit also knows the true binary label and can derive
per-record loss; its frozen primary attacker is a loss threshold.
Predeclared threshold, logistic-regression, and random-forest attackers
provide secondary sensitivity analysis. Target evaluation uses 25,194
balanced records per model.

MIA advantage is $\max_t(\mathrm{TPR}(t)-\mathrm{FPR}(t))$. We
additionally report ROC-AUC, balanced accuracy at the shadow-calibrated
threshold, and TPR at FPR ceilings 0.01 and 0.05. The IDS threshold
never transforms the raw probabilities supplied to an MIA.

\hypertarget{statistical-analysis}{%
\subsubsection{3.6 Statistical analysis}\label{statistical-analysis}}

For NSL-KDD, condition summaries are means and standard deviations
across the five target-training seeds. Two-sided t intervals use the
target seed as the uncertainty unit (four degrees of freedom).
DP-minus-non-private differences are paired by seed. MIA bootstraps use
1,000 replicates on identical target records where declared. These
intervals are conditional on the frozen dataset, split, architecture,
hyperparameters, and attackers. They do not represent uncertainty over
all possible datasets, architectures, or adversaries.

For UNSW-NB15, 95\% bootstrap intervals quantify record-sampling
uncertainty for one trained target. They do not capture
target-training-seed variation. We therefore do not test cross-dataset
equality or pool the two datasets.

\hypertarget{results}{%
\subsection{4. Results}\label{results}}

\hypertarget{nsl-kdd-utility}{%
\subsubsection{4.1 NSL-KDD utility}\label{nsl-kdd-utility}}

\begin{longtable}[]{@{}
  >{\raggedright\arraybackslash}p{(\columnwidth - 12\tabcolsep) * \real{0.1111}}
  >{\raggedleft\arraybackslash}p{(\columnwidth - 12\tabcolsep) * \real{0.1481}}
  >{\raggedleft\arraybackslash}p{(\columnwidth - 12\tabcolsep) * \real{0.1481}}
  >{\raggedleft\arraybackslash}p{(\columnwidth - 12\tabcolsep) * \real{0.1481}}
  >{\raggedleft\arraybackslash}p{(\columnwidth - 12\tabcolsep) * \real{0.1481}}
  >{\raggedleft\arraybackslash}p{(\columnwidth - 12\tabcolsep) * \real{0.1481}}
  >{\raggedleft\arraybackslash}p{(\columnwidth - 12\tabcolsep) * \real{0.1481}}@{}}
\toprule\noalign{}
\begin{minipage}[b]{\linewidth}\raggedright
Condition
\end{minipage} & \begin{minipage}[b]{\linewidth}\raggedleft
Actual epsilon
\end{minipage} & \begin{minipage}[b]{\linewidth}\raggedleft
Recall
\end{minipage} & \begin{minipage}[b]{\linewidth}\raggedleft
FNR
\end{minipage} & \begin{minipage}[b]{\linewidth}\raggedleft
F1
\end{minipage} & \begin{minipage}[b]{\linewidth}\raggedleft
FPR
\end{minipage} & \begin{minipage}[b]{\linewidth}\raggedleft
Average precision
\end{minipage} \\
\midrule\noalign{}
\endhead
\bottomrule\noalign{}
\endlastfoot
Non-private & --- & 0.7093 ± 0.0222 & 0.2907 ± 0.0222 & 0.8097 ± 0.0161
& 0.0561 ± 0.0227 & 0.9375 ± 0.0079 \\
DP-SGD epsilon about 4 & 3.9983 & 0.7134 ± 0.0138 & 0.2866 ± 0.0138 &
0.8017 ± 0.0082 & 0.0876 ± 0.0042 & 0.8907 ± 0.0063 \\
DP-SGD epsilon about 2 & 1.9990 & 0.7011 ± 0.0116 & 0.2989 ± 0.0116 &
0.7946 ± 0.0074 & 0.0838 ± 0.0015 & 0.8906 ± 0.0063 \\
\end{longtable}

Values are mean ± standard deviation over five target seeds at
validation-selected thresholds.

The epsilon-about-4 Recall difference from non-private is +0.00404 with
a 95\% paired t interval {[}-0.01334, 0.02142{]}; its F1 difference is
-0.00809 {[}-0.02241, 0.00622{]}. Neither establishes an improvement or
degradation across seeds. The operational costs are clearer: FPR
increases by 0.03145 {[}0.00460, 0.05830{]}, and average precision
decreases by 0.04682 {[}-0.05462, -0.03903{]}.

At epsilon about 2, Recall differs by -0.00826 {[}-0.03135, 0.01483{]}
and F1 by -0.01514 {[}-0.03364, 0.00335{]}, again uncertain. FPR
increases by 0.02772 {[}0.00038, 0.05507{]}, while average precision
decreases by 0.04693 {[}-0.05312, -0.04074{]}. Thus, neither private
condition is a confirmed Recall optimum, and both produce a worse
false-alarm/ranking profile under this protocol.

\begin{figure}
\centering
\researchfigure{../results/final_analysis/figures/10_false_alarms_and_average_precision.png}
\caption{NSL-KDD false-positive rate and average precision}
\end{figure}

\emph{Figure 1. Validation-selected F2 operating points over five
NSL-KDD target seeds. Additional false-positive and average-precision
diagnostics prevent selection by Recall alone.}

\hypertarget{nsl-kdd-membership-inference}{%
\subsubsection{4.2 NSL-KDD membership
inference}\label{nsl-kdd-membership-inference}}

\begin{longtable}[]{@{}
  >{\raggedright\arraybackslash}p{(\columnwidth - 8\tabcolsep) * \real{0.1579}}
  >{\raggedleft\arraybackslash}p{(\columnwidth - 8\tabcolsep) * \real{0.2105}}
  >{\raggedleft\arraybackslash}p{(\columnwidth - 8\tabcolsep) * \real{0.2105}}
  >{\raggedleft\arraybackslash}p{(\columnwidth - 8\tabcolsep) * \real{0.2105}}
  >{\raggedleft\arraybackslash}p{(\columnwidth - 8\tabcolsep) * \real{0.2105}}@{}}
\toprule\noalign{}
\begin{minipage}[b]{\linewidth}\raggedright
Condition
\end{minipage} & \begin{minipage}[b]{\linewidth}\raggedleft
Score-only AUC
\end{minipage} & \begin{minipage}[b]{\linewidth}\raggedleft
Label-aware AUC
\end{minipage} & \begin{minipage}[b]{\linewidth}\raggedleft
Score-only advantage
\end{minipage} & \begin{minipage}[b]{\linewidth}\raggedleft
Label-aware advantage
\end{minipage} \\
\midrule\noalign{}
\endhead
\bottomrule\noalign{}
\endlastfoot
Non-private & 0.50230 & 0.50088 & 0.00941 & 0.00786 \\
DP-SGD epsilon about 4 & 0.50301 & 0.50159 & 0.01107 & 0.01097 \\
DP-SGD epsilon about 2 & 0.50327 & 0.50146 & 0.01145 & 0.01130 \\
\end{longtable}

All values are close to chance in magnitude. Some conditional
across-seed AUC intervals exclude 0.5 by a few thousandths, so it would
be inaccurate to say that every NSL-KDD interval contains chance. The
more relevant leakage-reduction question is paired. For epsilon about 4,
paired AUC differences are +0.00071 {[}-0.00111, 0.00254{]} score-only
and +0.00071 {[}-0.00094, 0.00236{]} label-aware. For epsilon about 2,
they are +0.00097 {[}-0.00076, 0.00270{]} and +0.00058 {[}-0.00126,
0.00242{]}. All primary and secondary paired AUC intervals cross zero.
No tested private condition shows a measurable AUC reduction.

The secondary epsilon-about-2 label-aware advantage increases by 0.00345
{[}0.00068, 0.00621{]}. This is a small exploratory result from a
secondary attacker and does not show a broad privacy degradation, but it
reinforces that the evidence cannot be summarized as leakage reduction.

\begin{figure}
\centering
\researchfigure{../results/final_analysis/figures/04_mia_auc_vs_epsilon.png}
\caption{NSL-KDD MIA AUC}
\end{figure}

\emph{Figure 2. Fixed primary attacks with conditional five-seed t
intervals. Near-chance magnitude is not proof that leakage is absent.}

\hypertarget{supplementary-unsw-nb15-check}{%
\subsubsection{4.3 Supplementary UNSW-NB15
check}\label{supplementary-unsw-nb15-check}}

\begin{longtable}[]{@{}
  >{\raggedright\arraybackslash}p{(\columnwidth - 12\tabcolsep) * \real{0.1111}}
  >{\raggedleft\arraybackslash}p{(\columnwidth - 12\tabcolsep) * \real{0.1481}}
  >{\raggedleft\arraybackslash}p{(\columnwidth - 12\tabcolsep) * \real{0.1481}}
  >{\raggedleft\arraybackslash}p{(\columnwidth - 12\tabcolsep) * \real{0.1481}}
  >{\raggedleft\arraybackslash}p{(\columnwidth - 12\tabcolsep) * \real{0.1481}}
  >{\raggedleft\arraybackslash}p{(\columnwidth - 12\tabcolsep) * \real{0.1481}}
  >{\raggedleft\arraybackslash}p{(\columnwidth - 12\tabcolsep) * \real{0.1481}}@{}}
\toprule\noalign{}
\begin{minipage}[b]{\linewidth}\raggedright
Condition
\end{minipage} & \begin{minipage}[b]{\linewidth}\raggedleft
Actual epsilon
\end{minipage} & \begin{minipage}[b]{\linewidth}\raggedleft
Threshold
\end{minipage} & \begin{minipage}[b]{\linewidth}\raggedleft
Recall
\end{minipage} & \begin{minipage}[b]{\linewidth}\raggedleft
F1
\end{minipage} & \begin{minipage}[b]{\linewidth}\raggedleft
FPR
\end{minipage} & \begin{minipage}[b]{\linewidth}\raggedleft
Average precision
\end{minipage} \\
\midrule\noalign{}
\endhead
\bottomrule\noalign{}
\endlastfoot
Non-private & --- & 0.24 & 0.998831 & 0.853572 & 0.418432 & 0.981609 \\
DP-SGD epsilon about 4 & 3.995481 & 0.02 & 0.999846 & 0.852478 &
0.423784 & 0.968490 \\
DP-SGD epsilon about 2 & 1.995670 & 0.02 & 0.999890 & 0.852436 &
0.424000 & 0.967663 \\
\end{longtable}

At the selected F2 operating point, both private models retain F1 within
0.0012 of non-private and slightly increase Recall. They also increase
FPR by about 0.0054--0.0056 and reduce average precision by about
0.0131--0.0139. More importantly, absolute FPR is approximately 42\% for
every condition. At the default 0.5 threshold it remains 27.3\% for
non-private and 40.7\% for the private models. Near-perfect Recall
therefore does not imply operational deployability.

Score-only MIA AUCs are 0.499363 {[}0.493715, 0.505200{]}, 0.500772
{[}0.494636, 0.506606{]}, and 0.500000 {[}0.500000, 0.500000{]}.
Label-aware AUCs are 0.504553 {[}0.498349, 0.510101{]}, 0.504064
{[}0.498273, 0.509983{]}, and 0.504033 {[}0.497777, 0.509767{]}. Every
interval contains chance. The epsilon-about-2 score-only attacker emits
constant scores; its zero advantage and separated paired advantage
difference are a degenerate measurement, not proof of absent leakage.

\begin{figure}
\centering
\researchfigure{../results/unsw_nb15_external_validation/figures/01_unsw_ids_utility.png}
\caption{UNSW-NB15 utility}
\end{figure}

\emph{Figure 3. Single-seed UNSW-NB15 utility at validation-selected F2
thresholds. The table and FPR must accompany the high Recall values.}

\begin{figure}
\centering
\researchfigure{../results/unsw_nb15_external_validation/figures/02_unsw_mia_auc.png}
\caption{UNSW-NB15 MIA AUC}
\end{figure}

\emph{Figure 4. Single-seed UNSW-NB15 MIA AUC with record-level
bootstrap intervals.}

\hypertarget{discussion}{%
\subsection{5. Discussion}\label{discussion}}

\hypertarget{answers-to-the-research-questions}{%
\subsubsection{5.1 Answers to the research
questions}\label{answers-to-the-research-questions}}

\textbf{RQ1: DP-SGD changes the utility profile more clearly than it
changes Recall.} Across five NSL-KDD seeds, the epsilon-about-4 Recall
gain is too uncertain to call an improvement, and the epsilon-about-2
Recall change is also uncertain. Both private conditions have materially
lower average precision and higher FPR. The result illustrates why
choosing a privacy budget from Recall or FNR alone is unsafe. The F2
threshold moves the operating point toward attacks, especially for
private models whose selected thresholds are much lower, but it cannot
restore ranking quality.

\textbf{RQ2: the evaluated attacks do not demonstrate DP-induced leakage
reduction.} The non-private model's primary AUCs are already
approximately 0.5. Private-model AUCs remain in the same narrow range,
and paired intervals cross zero. This is an empirical floor effect: an
attack cannot show a large reduction from a baseline it already fails to
distinguish. The result does not conflict with formal DP. The DP
guarantee is worst-case and mechanism-based; the MIA is a finite,
attack-relative measurement. Conversely, a weak MIA cannot be used to
claim privacy for the non-private model.

\textbf{RQ3: UNSW-NB15 is qualitatively consistent only at a broad
level.} The supplementary study again shows close selected-threshold
F1/Recall, slightly worse private-model FPR and average precision, and
near-chance overall MIA. However, it has one target seed, different data
and preprocessing, and an exceptionally high absolute FPR. It supports a
bounded consistency statement, not a universal generalization or formal
comparison between datasets.

\hypertarget{what-can-be-selected-from-the-tested-conditions}{%
\subsubsection{5.2 What can be selected from the tested
conditions?}\label{what-can-be-selected-from-the-tested-conditions}}

No tested epsilon is a universal optimum. Epsilon about 4 provides the
highest mean NSL-KDD Recall, but the paired interval includes both
benefit and harm; its higher FPR and lower average precision are
clearer. Epsilon about 2 gives a tighter formal bound but slightly lower
mean F1 and does not improve measured MIA. A deployment decision would
need an explicit privacy valuation and a false-alarm cost, neither of
which the benchmark datasets provide. We therefore present a tested
frontier rather than naming a winner.

\hypertarget{implications-for-empirical-privacy-evaluation}{%
\subsubsection{5.3 Implications for empirical privacy
evaluation}\label{implications-for-empirical-privacy-evaluation}}

The negative MIA result is scientifically useful. It demonstrates that
adding a formal guarantee does not ensure that a population-average
audit will measure a difference. Stronger attacks may reveal leakage
that the present features and shadows miss, especially for individual
vulnerable records \citep{carlini2022membership, aerni2024evaluations}.
A future privacy audit should be motivated by a specific claim:
LiRA/RMIA for stronger population inference, canaries for an empirical
lower bound on DP, a label-only attack for restricted APIs, or an
adaptive attack that knows the defense. Running more attacks solely
after observing null results would introduce selection bias.

\hypertarget{operational-ids-implications}{%
\subsubsection{5.4 Operational IDS
implications}\label{operational-ids-implications}}

F2 selection encodes a preference for Recall, not a universal
operational objective. On NSL-KDD, the private models trade a modest
threshold-specific Recall profile for roughly three additional false
positives per hundred normal records and a substantial decrease in
average precision. On UNSW-NB15, the selected points flag more than four
in ten normal test records. These figures are not deployment claims;
they are evidence that a single aggregate score can hide an unacceptable
alert burden. Future deployment studies should select thresholds from
explicit alert capacity and cost, then validate under temporal and site
shift.

\hypertarget{limitations-and-threats-to-validity}{%
\subsection{6. Limitations and threats to
validity}\label{limitations-and-threats-to-validity}}

\begin{enumerate}
\def\labelenumi{\arabic{enumi}.}
\tightlist
\item
  \textbf{Formal scope.} DP covers optimization conditional on fixed
  preprocessing. Categorical vocabularies, scaling statistics,
  architecture, hyperparameters, and data selection are outside the
  stated guarantee.
\item
  \textbf{Research randomness.} \texttt{secure\_mode=False} is recorded.
  The implementation is suitable for experimental comparison, not a
  production cryptographic claim.
\item
  \textbf{Attack coverage.} The primary attacks are score-only logistic
  regression and a label-aware loss threshold. LiRA, RMIA,
  canary/worst-case auditing, label-only attacks, white-box access,
  adaptive attacks, and auxiliary-information variation are absent.
\item
  \textbf{Baseline floor.} Near-chance non-private MIA leaves little
  measurable headroom. Failure to distinguish members under these
  attacks is not evidence that the non-private model is safe.
\item
  \textbf{Dataset age and realism.} NSL-KDD is an older derived
  benchmark. UNSW-NB15 is more recent but still a benchmark generated in
  a controlled environment. Neither establishes performance on live
  organizational traffic.
\item
  \textbf{External-validation power.} UNSW-NB15 has one target-training
  seed. Its bootstrap intervals condition on that fitted model and do
  not measure training variation.
\item
  \textbf{Model and task scope.} The primary claim concerns one MLP
  family and binary classification. It does not cover multiclass
  attribution, other architectures, online learning, federated learning,
  or other DP mechanisms.
\item
  \textbf{Operating policy.} F2 thresholding emphasizes Recall.
  Different costs or alert budgets could reorder the practical
  preference among models.
\item
  \textbf{Conditional inference.} Five-seed t intervals are based on a
  small number of seeds and one locked split. Subgroup and
  epsilon-about-8 findings are exploratory/single-seed context.
\item
  \textbf{No adversarial-evasion evaluation.} DP and adversarial
  robustness can interact, but the present experiments do not test
  protocol-valid evasion or poisoning.
\end{enumerate}

\hypertarget{reproducibility-and-ethics}{%
\subsection{7. Reproducibility and
ethics}\label{reproducibility-and-ethics}}

The repository records dataset hashes, split sizes, preprocessing
manifests, model and attacker configurations, privacy-accounting
parameters, warning logs, per-seed summaries, paired analyses, figures,
and evidence manifests. Large model states and per-record score exports
remain in the canonical external evidence bundles; compact aggregate
evidence and checksums are committed. Experiment 08 performs analysis
only and starts no training, threshold tuning, or attacker fitting.
Experiment 09 records all 18 target/shadow completions and verifies
official dataset identities.

This work uses public benchmark network-traffic datasets and does not
collect new human-subject data. Nevertheless, membership inference is
dual-use. The study reports aggregate attacks for defensive evaluation,
avoids releasing private real-world traffic, and does not claim that a
weak attack makes a model safe to deploy.

\hypertarget{conclusion}{%
\subsection{8. Conclusion}\label{conclusion}}

We evaluated non-private and formally accounted DP-SGD training for a
binary MLP IDS using a five-seed NSL-KDD study and a single-seed
UNSW-NB15 check. DP-SGD at epsilon about 4 and 2 retained broadly
similar selected-threshold Recall/F1, but increased the false-positive
burden and reduced average precision. Population-level score-only and
label-aware MIAs remained near chance, and paired AUC comparisons did
not show leakage reduction. The correct conclusion is therefore bounded:
the formal DP mechanism supplies privacy evidence conditional on fixed
preprocessing, while the tested MIAs reveal no measurable improvement
over an already near-chance baseline. Privacy budgets should not be
chosen from Recall alone, and empirical attack failure should not be
confused with privacy certification.

\bibliography{references}

\end{document}
